\begin{helper}[A maximal fixed-Gamma augmentation chain]
Let $J$ carry a hierarchy $K_J$ and a good witness $\mathcal G_J$ under a
coloring $\sigma$.  Put $d=K_J.\mathit{dividerCount}$,
$A=K_J.\mathit{starred}(d)$, $D=D^{K_J}_{n,d}$, and
$\Gamma_\eta=\mathcal G_J.\Gamma^d_\eta$.  For an ordered tree $W$, let
$C_W(\eta)$ be the colors on displayed edges whose positions are at least
$|A|$ and strictly below both $|W|$ and the totalized first position in
$W$ missing $\eta$.

There are $m$, ordered trees $U_0=J,\ldots,U_m=U$, and a terminal tree
$U$ such that each successor is one
$\sigma$-\noderef{TashkinovAugmentation}, satisfies
$\Gamma_\eta\cap C_{U_{r+1}}(\eta)=\varnothing$ for every $\eta\in D$,
has length $|J|+r+1$, and properly enlarges the vertex set.  Every earlier
$U_r$ is an exact ordered initial segment of every later $U_s$, the whole
chain is a \noderef{TashkinovAugmentationSequence} from $J$ to $U$, and
the endpoint satisfies the same avoidance conditions.  Moreover, $U$ has
no one-step Tashkinov augmentation satisfying all of those fixed-Gamma
avoidance conditions.
\end{helper}
\begin{proof}
Let $d,A,D$, and $\Gamma_\eta$ be as in the statement.  By
\noderef{HierarchySegmentColors}, $C_J(\eta)$ is the last-index segment
color set in the avoidance field of \noderef{GoodHierarchy}.  Therefore
that field gives
$\Gamma_\eta\cap C_J(\eta)=\varnothing$ for every $\eta\in D$.

Call a finite chain admissible when it starts at $J$, every successor is a
one-step $\sigma$-\noderef{TashkinovAugmentation}, and every successor
satisfies all the displayed fixed-Gamma avoidance conditions.  The
length-zero chain is admissible by the preceding paragraph.  Every
successor has length one more than its predecessor and contains it as an
exact \noderef{TreeSequenceInitialSegment}, by the definition of
Tashkinov augmentation.  Hence all old vertex positions have the same
entries, and the successor has exactly one additional vertex position.
Injectivity in \noderef{TreeSequence} makes that entry different from
every old entry.  Thus the successor properly enlarges
$\noderef{TreeSequenceVertexSet}$.

After $q$ successors, the vertex set therefore contains $q$ distinct
vertices in addition to $V(J)$.  Since $V$ is finite, $q\leq
|V|-|V(J)|$.  The nonempty set of attainable chain lengths is consequently
a subset of a finite interval of natural numbers, so it has a greatest
member $m$.  Choose an admissible chain $U_0,\ldots,U_m$ of that length and
put $U=U_m$.  An admissible one-step successor of $U$ could be appended to
this chain, producing length $m+1$ and contradicting maximality.  This is
exactly the asserted endpoint maximality.  Endpoint avoidance holds at
$m=0$ by the initial good-hierarchy avoidance and, at $m>0$, by the
admissibility condition on the last successor.

Composition of the entrywise equalities in
\noderef{TreeSequenceInitialSegment} proves by induction on $s-r$ that
$U_r$ is an exact ordered initial segment of $U_s$ whenever $r\leq s$.
The one-more length equations give $|U_r|=|J|+r$ by induction.  Finally,
starting with the reflexive case and adjoining each one-step relation with
the transitive-tail constructor proves, from the definition of
\noderef{TashkinovAugmentationSequence}, that the complete chain is such a
sequence from $J$ to $U$.  All claimed chain properties follow.
\end{proof}
